% Capitulo: Na prática: projetos de engenharia de dados
\chapter{Na prática: projetos de engenharia de dados}
\label{cap:proj-engenharia}

\rascunho{
  \textbf{Ideia do capítulo.} Como o modelo funciona em pipelines, integrações e plataforma.
  \begin{itemize}
    \item Épicos e histórias técnicas: como escrever história quando o usuário é outro sistema.
    \item Débito técnico e onde ele entra no quadro.
    \item Projeto contra sustentação: duas raias ou dois times?
    \item Projetos paralelos puxados por um projeto de negócio, como o de master data
      (\veja{sec:master-data-puxado}): épico próprio, histórias amarradas às do projeto
      principal.
    \item Relato de um projeto real, sem nomes de empresa ou cliente.
  \end{itemize}
}

\section{Histórias quando o usuário é um sistema}
\label{sec:historias-tecnicas}

\pendente{escrever}

\section{Débito técnico}
\label{sec:debito-tecnico}

\pendente{escrever}

\section{Projeto e sustentação}
\label{sec:projeto-sustentacao}

\pendente{escrever}
